\section{Resumen de esfuerzo}
\label{anexo:resumen-esfuerzo}

El esfuerzo de F.A.M.A. se expresa en horas-persona, que corresponden a
la suma de las horas de trabajo de los integrantes del equipo. Según las
cifras consolidadas en el capítulo~9, la planificación inicial consideró
320 horas-persona y el esfuerzo real fue de 317 horas-persona. La
Tabla~\ref{tab:esfuerzo-fama} presenta el desglose disponible del
esfuerzo real y la participación de cada actividad en el total.

\begin{table}[!htbp]
    \centering
    \caption{Esfuerzo planificado y real del equipo de F.A.M.A.}
    \label{tab:esfuerzo-fama}
    \small
    \setlength{\tabcolsep}{4pt}
    \renewcommand{\arraystretch}{1.15}
    \begin{tabular}{@{}p{0.47\linewidth}rrr@{}}
        \hline
        \textbf{Actividad} & \shortstack[r]{\textbf{Planificado}\\(h-p)} &
        \shortstack[r]{\textbf{Real}\\(h-p)} & \textbf{\% del real} \\
        \hline
        Desarrollo de software (programación, pruebas automatizadas y control de versiones)
        & --- & 132 & 41,64\% \\
        Elaboración de la memoria (redacción y corrección)
        & --- & 185 & 58,36\% \\
        \hline
        \textbf{Total del equipo} & \textbf{320} & \textbf{317} & \textbf{100\%} \\
        \hline
    \end{tabular}
    \par\smallskip
    \parbox{0.97\linewidth}{\footnotesize
    h-p significa horas-persona. El guion indica que la memoria
    consigna la planificación únicamente como total del equipo.
    El registro disponible agrupa el esfuerzo real por actividad,
    sin separar horas por integrante o por tarea técnica.}
\end{table}

\subsection{Comparación con la planificación}

El esfuerzo real quedó 3 horas-persona por debajo del total
planificado, una desviación de $-0{,}94$\% respecto de las 320
horas-persona previstas. Las conclusiones del proyecto señalan que el
procesamiento tensorial, la estabilización del uso de memoria de la
GPU y el manejo de concurrencia requirieron más dedicación técnica de
la prevista. El desglose disponible permite observar la distribución
entre desarrollo y elaboración de la memoria, mientras que la
variación dentro de las tareas técnicas se describe cualitativamente.

La carta Gantt del apartado~\ref{anexo:gantt} sitúa las actividades en
el calendario y muestra los cambios en sus plazos. Este resumen
cuantifica la dedicación acumulada del equipo. Por su parte, el
apartado~\ref{anexo:estimacion-cu} aplica una tasa de referencia a los
puntos de casos de uso para obtener un esfuerzo teórico de desarrollo.
El presupuesto de 320 horas-persona corresponde a la planificación
académica del equipo e incluye la elaboración de la memoria.

% Fuente de las horas planificadas y reales: capítulo 9 de la memoria.
% Participaciones: 132/317 y 185/317. Desviación: (317-320)/320.
